% !PRINTED 604
\weberpage{604}

This question has found a quite general answer in a memoir by \gesperrt{Hilbert}\footnote{Ueber die Irreducibilit\"at ganzer rationaler Functionen mit ganzzahligen Co\"efficienten. Journal f\"ur Mathematik, Bd. 110.}, in which the general theorem is proved that in an irreducible function of arbitrarily many variables one can substitute for an arbitrary part of the variables rational numbers such that an irreducible function of the remaining variables results. We cannot enter upon this general theorem here.

We shall, however, answer the question posed much more simply, though by far not so generally, by showing that equations without affect can be found for every prime degree.

We have proved in \S.\,153, 9. that a transitive permutation group of $n$ figures which is not the symmetric group, if $n$ is a prime, can contain no single transposition.

By an equation with an affect we have understood one whose group is not the symmetric group. If, then, the irreducible equation $f(x) = 0$ has an affect and its degree is a prime, its group $P$ must be transitive, and it can contain no transposition of two roots. If therefore we adjoin any $n - 2$ of the roots to the domain of rationality, the group must reduce to the unit group, since besides this only the permutation of the two last roots could still remain, which does not occur in $P$. From this it follows, then, that the two last roots are contained in the extended field $\Omega$, or the theorem:

1. \gesperrt{If an irreducible equation whose degree $n$ is a prime has an affect, then any two of its roots can be expressed rationally in terms of the remaining ones.}

From this follows as a corollary:

2. \gesperrt{If the field $\Omega$ contains only real numbers, then an irreducible equation of prime degree $n$ with an affect in $\Omega$ cannot have two imaginary and $n - 2$ real roots.}

Now there are, however, innumerable equations of every arbi-
%CONT>>
